\chapter{자원: 회로는 공간을 차지한다}\label{ch:resources}

프로그램의 비용은 주로 실행 시간이다. 회로의 비용은 둘이다. 부품이 차지하는 \emph{자원}과 신호가 지나는
데 걸리는 \emph{시간}. 이 장은 FPGA 안에 어떤 부품이 있고, 코드가 그 부품을 얼마나 쓰는지 가늠하는 법을
설명한다. 이 감각이 없으면 \ref{ch:lesson}장의 사고가 난다.

\section{FPGA 안에 있는 것}

\Why FPGA는 만들어진 뒤에 회로를 바꿀 수 있는 칩이다. 게이트를 그때그때 새로 만들 수는 없으므로,
FPGA는 ``어떤 논리 함수로든 설정할 수 있는 작은 부품''을 수만 개에서 수십만 개 깔아 두고 그 사이를
설정 가능한 배선으로 잇는다. 그 작은 부품이 \term{LUT}{look-up table, 조회표}다.

LUT는 이름 그대로 표다. 입력이 6개인 LUT(LUT6)는 $2^6 = 64$칸짜리 1비트 메모리이고, 입력 6비트를 주소로
삼아 그 칸의 값을 내놓는다. 64칸에 무엇을 적어 두느냐에 따라 입력 6개의 \emph{어떤} 논리 함수든 된다
(그림~\ref{fig:g-lut}). 합성 도구가 하는 일은 설계자가 적은 논리를 6입력 이하의 조각으로 쪼개 LUT에
하나씩 담는 것이다.

\begin{figure}[ht]
\centering
\begin{tikzpicture}[x=1cm,y=1cm]
  % 진리표
  \node[font=\small\bfseries] at (1.5,3.45) {LUT에 적힌 표 (일부)};
  \matrix[matrix of nodes,nodes={font=\ttfamily\footnotesize,minimum width=5mm,minimum height=4.2mm,inner sep=0pt},
          column sep=0pt,row sep=0pt] (m) at (1.5,0.9) {
     a & b & c & |[text=cgray]|\dots & {} & y \\
     0 & 0 & 0 & {} & {} & 0 \\
     0 & 0 & 1 & {} & {} & 1 \\
     0 & 1 & 0 & {} & {} & 0 \\
     0 & 1 & 1 & {} & {} & 1 \\
     1 & 0 & 0 & {} & {} & 0 \\
     1 & 0 & 1 & {} & {} & 1 \\
     1 & 1 & 0 & {} & {} & 1 \\
     1 & 1 & 1 & {} & {} & 1 \\
  };
  \draw (m-1-1.south west) -- (m-1-6.south east);
  \draw (m-1-5.north east) -- (m-9-5.south east);
  \node[note] at (1.5,-1.85) {\code{y = (a \& b) | c}};
  % LUT 기호
  \begin{scope}[xshift=6.6cm]
  \node[font=\small\bfseries] at (1.6,2.9) {LUT6 하나};
  \node[draw,fill=cbluebg,draw=cblue,minimum width=22mm,minimum height=30mm,align=center] (lut) at (1.6,0.8) {64칸\\메모리};
  \foreach \i/\l in {0/a,1/b,2/c,3/{},4/{},5/{}} {
    \draw[<-] ([yshift=11mm-\i*4.4mm]lut.west) -- ++(-0.8,0) node[left,font=\ttfamily\footnotesize] {\l};
  }
  \draw[->] (lut.east) -- ++(0.8,0) node[right,font=\ttfamily\footnotesize] {y};
  \node[note] at (1.6,-1.4) {입력 6비트가 주소, 출력은 그 칸의 값};
  \end{scope}
  % FF
  \begin{scope}[xshift=11.6cm]
  \node[font=\small\bfseries] at (1.2,2.9) {짝을 이루는 플립플롭};
  \node[reg,minimum width=12mm,minimum height=14mm] (ff) at (1.2,0.8) {FF};
  \pic at (ff.south) {clkmark};
  \draw[<-] (ff.west) -- ++(-0.6,0);
  \draw[->] (ff.east) -- ++(0.6,0);
  \node[note] at (1.2,-1.4) {1비트 기억};
  \end{scope}
\end{tikzpicture}
\caption{LUT는 작은 메모리로 만든 만능 게이트다. 그림~\ref{fig:g-assign}의 게이트 두 개는 LUT 하나에
들어간다(입력 세 개만 쓴다).}
\label{fig:g-lut}
\end{figure}

\General \code{y = (a \& b) | c}는 입력이 셋이므로 LUT 하나다. 게이트가 몇 개인지는 중요하지 않다.
중요한 것은 \emph{출력 하나가 입력 몇 개에 의존하는가}다. 입력이 여섯 이하이면 아무리 복잡한 식도
LUT 하나이고, 일곱 이상이면 LUT 여러 개를 엮어야 한다.

FPGA에는 LUT 말고도 몇 가지 부품이 있다(표~\ref{tab:g-res}).

\begin{table}[ht]
\centering
\caption{FPGA의 주요 자원. 수량은 이 프로젝트의 대상 칩 두 종의 대략적인 값이다.}
\label{tab:g-res}
\small
\begin{tabularx}{\linewidth}{l>{\raggedright\arraybackslash}X>{\raggedright\arraybackslash}p{27mm}>{\raggedright\arraybackslash}p{24mm}}
\toprule
자원 & 하는 일 & XC7A200T (Nexys Video) & XC7Z020 (PYNQ) \\
\midrule
LUT & 6입력 이하의 임의의 논리 함수. 조합 논리가 여기에 담긴다 & 약 13만 4천 & 약 5만 3천 \\
플립플롭(FF) & 1비트 기억. 레지스터가 여기에 담긴다 & LUT마다 딸려 있다 & \\
분산 RAM (LUT RAM) & LUT의 64칸 메모리를 그대로 작은 RAM으로 쓰는 것. 작고 빠른 표에 알맞다 & LUT의 일부를 전용 & \\
블록 RAM & 큰 메모리 전용 블록(수만 비트 단위). 읽는 데 한 클럭이 든다 & 365개 & 140개 \\
DSP 블록 & 곱셈기. 신호 처리에 쓰인다 & (이 프로젝트에서는 쓰지 않는다) & \\
\bottomrule
\end{tabularx}
\end{table}

\Ours 봄베의 회로에는 곱셈이 없고 큰 메모리도 필요 없다. 자원의 거의 전부가 LUT(전파 논리와 치환 계산)와
플립플롭(676비트 상태와 치환표)이다. 그래서 이 프로젝트에서 ``회로가 큰가''는 곧 ``LUT를 몇 개 쓰는가''다.

\section{선택 회로는 비싸다}\label{sec:mux-cost}

회로에서 LUT를 가장 많이 잡아먹는 것은 덧셈도 비교도 아니고 \term{선택 회로}{multiplexer, mux}다.
여러 입력 가운데 하나를 골라 내보내는 회로다.

\Why 선택 회로의 출력은 \emph{모든} 데이터 입력과 선택 신호에 의존한다. 넷 중 하나를 고르는 4:1 선택
회로는 데이터 4비트와 선택 2비트, 모두 여섯 입력이므로 LUT6 하나에 꼭 맞는다. 그보다 큰 선택 회로는
4:1을 층층이 쌓아야 한다. 16:1은 4:1 네 개로 16개를 4개로 줄이고 다시 4:1 하나로 고르므로 LUT 다섯 개다.
이런 식으로 $N$개 가운데 하나를 고르려면 LUT가 대략 $N/3$개 든다($N/4 + N/16 + \cdots$). 그리고 이것은
\emph{출력 1비트당}의 값이다(그림~\ref{fig:g-muxtree}).

\begin{figure}[ht]
\centering
\begin{tikzpicture}[x=1cm,y=1cm]
  \node[font=\small\bfseries] at (0.9,2.6) {4:1 = LUT 1개};
  \node[mux,minimum width=20mm] (m0) at (0.9,0.9) {4:1};
  \foreach \i in {0,1,2,3} { \draw[<-] ([yshift=7.5mm-\i*5mm]m0.west) -- ++(-0.6,0); }
  \draw[->] (m0.east) -- ++(0.5,0);
  \draw[<-] (m0.south) -- ++(0,-0.45) node[below,font=\footnotesize] {선택 2비트};
  \begin{scope}[xshift=5.2cm]
  \node[font=\small\bfseries] at (2.0,2.6) {16:1 = LUT 5개};
  \foreach \k/\yy in {0/2.0,1/1.0,2/0.0,3/-1.0} {
    \node[mux,minimum width=9mm,minimum height=5mm,font=\scriptsize] (a\k) at (0.8,\yy) {4:1};
    \foreach \i in {0,1,2,3} { \draw ([yshift=3mm-\i*2mm]a\k.west) -- ++(-0.35,0); }
  }
  \node[mux,minimum width=20mm] (b) at (3.4,0.5) {4:1};
  \foreach \k in {0,1,2,3} { \draw[->] (a\k.east) -- ++(0.5,0) |- ([yshift=7.5mm-\k*5mm]b.west); }
  \draw[->] (b.east) -- ++(0.5,0);
  \end{scope}
  \begin{scope}[xshift=11.3cm]
  \node[font=\small\bfseries] at (1.6,2.6) {크기의 어림 (출력 1비트당)};
  \node[font=\small,align=left,anchor=north west] at (-0.4,2.2) {
     \begin{tabular}{@{}rl@{}}
       4:1 & LUT 1개 \\
       16:1 & 5개 \\
       26:1 & 열 개 안팎 \\
       676:1 & 이백 개 이상 \\
     \end{tabular}};
  \end{scope}
\end{tikzpicture}
\caption{선택 회로의 크기. 고를 대상이 많아지면 그에 비례해 커진다.}
\label{fig:g-muxtree}
\end{figure}

\General 선택 회로는 코드에서 눈에 잘 띄지 않는다. 다음 두 줄은 모두 선택 회로다. 특히 둘째 줄, 변수로
배열의 원소를 고르는 것은 프로그램에서는 가장 값싼 연산이지만 회로에서는 ``배열의 모든 원소를 선택 회로의
입력에 연결한다''는 뜻이다.

\begin{lstlisting}[language=SV,numbers=none]
assign y = sel ? a : b;      // 2:1 mux
assign z = table[index];     // (size of table):1 mux, for every bit of z
\end{lstlisting}

\Ours 봄베의 회로는 선택 회로투성이다. ``26글자 가운데 하나''를 고르는 일이 곳곳에 있기 때문이다.
로터를 통과하는 것은 배선표 26칸에서 한 칸을 고르는 일이고(5비트 폭의 26:1), 전파 논리에서 ``케이블 $a$의
$y$번 선''을 읽는 것도 선택이다. 코드~\ref{lst:g-spread}의 \code{live\_q[cable\_a + y]}가 그것인데,
여기서 고르는 대상은 26개가 아니라 676개다. 이 한 줄이 이 프로젝트의 가장 큰 사고의 원인이었다.
\ref{ch:lesson}장에서 다룬다.

\section{표는 메모리에 담는다}

\Why 같은 계산을 되풀이하는 것보다 계산해 둔 결과를 표에서 읽는 편이 싸다. 프로그램에서도 쓰는 방법이지만
회로에서는 의미가 더 크다. 계산 회로는 쓰든 안 쓰든 자리를 차지하고 있는 반면, 메모리는 같은 양의 정보를
훨씬 좁은 자리에 담는다. FPGA에는 메모리가 두 종류 있다. LUT를 메모리로 쓰는 \term{분산 RAM}{distributed
RAM}은 작지만 주소를 주면 같은 클럭 안에 값이 나오고, \term{블록 RAM}{block RAM}은 크지만 읽는 데 한 클럭이
걸린다.

\General 소리를 내는 회로가 사인파를 필요로 한다고 하자. 사인 함수를 계산하는 회로를 만드는 대신 한 주기의
값을 256칸짜리 표에 넣어 두고 차례로 읽어 낸다. 표는 블록 RAM 하나의 일부에 들어간다.

\Ours 로터의 배선표는 바뀌지 않는 상수 표다. 기준 RTL에서는 \code{wirings\_pkg}에 상수로 들어 있고
로터를 통과하는 함수가 그 표를 읽는다. FPGA 엔진은 한 걸음 더 나아가, 드럼 위치마다 다시 계산하던
스크램블러의 치환표를 미리 계산해 분산 RAM에 넣어 두고 읽어 쓴다. 다음 절의 주제다.

\section{자원과 시간을 맞바꾼다}\label{sec:tradeoff}

회로 설계에서 가장 자주 하는 결정은 ``부품을 더 쓰고 빨리 끝낼 것인가, 부품을 아끼고 오래 걸릴 것인가''다.

\Why 같은 계산을 하는 회로는 하나가 아니다. 계산에 필요한 연산을 전부 따로따로 부품으로 놓으면 한 클럭에
끝나지만 부품이 많이 든다. 부품 하나를 여러 클럭에 걸쳐 돌려 쓰면 부품은 적게 들지만 클럭이 여러 번
필요하다. 프로그램에서는 이 선택이 없다. 처리 장치가 하나이므로 언제나 뒤쪽이다.

\General 수 여덟 개의 합을 구한다고 하자(그림~\ref{fig:g-tradeoff}). 왼쪽은 덧셈기 일곱 개를 나무 모양으로
엮어 한 클럭에 답을 낸다. 오른쪽은 덧셈기 하나와 누산 레지스터 하나로 여덟 클럭에 걸쳐 하나씩 더한다.
신호 처리의 필터도 같다. 탭이 여덟 개인 필터는 곱셈기 여덟 개로 한 클럭에 계산할 수도 있고, 곱셈기 하나로
여덟 클럭에 걸쳐 계산할 수도 있다. 표본이 들어오는 속도가 느리면 뒤쪽으로 충분하다.

\begin{figure}[ht]
\centering
\begin{tikzpicture}[x=1cm,y=1cm,
   add/.style={draw,circle,fill=cbluebg,draw=cblue,minimum size=6mm,inner sep=0pt,font=\small}]
  \node[font=\small\bfseries] at (2.6,3.4) {병렬: 덧셈기 7개, 1클럭};
  \foreach \i in {0,...,7} { \node[font=\ttfamily\scriptsize] (x\i) at (\i*0.75,2.6) {x\i}; }
  \foreach \k/\a/\b in {0/0/1,1/2/3,2/4/5,3/6/7} {
    \node[add] (s\k) at (\k*1.5+0.375,1.7) {+};
    \draw[->] (x\a) -- (s\k); \draw[->] (x\b) -- (s\k);
  }
  \node[add] (t0) at (1.125,0.8) {+}; \node[add] (t1) at (4.125,0.8) {+};
  \draw[->] (s0) -- (t0); \draw[->] (s1) -- (t0); \draw[->] (s2) -- (t1); \draw[->] (s3) -- (t1);
  \node[add] (u) at (2.625,-0.1) {+};
  \draw[->] (t0) -- (u); \draw[->] (t1) -- (u);
  \draw[->] (u) -- ++(0,-0.6) node[below,font=\small] {합};
  \begin{scope}[xshift=9.0cm]
  \node[font=\small\bfseries] at (2.0,3.4) {직렬: 덧셈기 1개, 8클럭};
  \node[font=\ttfamily\small] (xin) at (-0.2,1.6) {x[i]};
  \node[add,minimum size=8mm] (p) at (1.3,1.6) {+};
  \node[reg,minimum width=14mm] (acc) at (3.3,1.6) {누산};
  \pic at (acc.south) {clkmark};
  \draw[->] (xin) -- (p); \draw[->] (p) -- (acc);
  \draw[->] (acc.east) -- ++(0.5,0) |- (2.3,0.5) -| (p.south);
  \draw[->] (acc.east) ++(0.5,0) -- ++(0.6,0) node[right,font=\small] {합};
  \node[note] at (2.0,-0.3) {클럭마다 수 하나를 더한다};
  \end{scope}
\end{tikzpicture}
\caption{같은 계산의 두 가지 구현.}
\label{fig:g-tradeoff}
\end{figure}

\Ours 봄베에서 이 선택은 스크램블러의 치환표를 언제, 어떻게 계산하는가에 나타난다.

실제 기계를 떠올려 보자. 봄베의 드럼 세 개 가운데 하나는 계속 돌지만 나머지 둘은 그 드럼이 한 바퀴 돌 때
한 칸씩만 움직인다. 26위치 동안 스크램블러의 ``안쪽''(가운데 로터, 왼쪽 로터, 반사판)은 그대로다.

\begin{itemize}
\item \textbf{기준 RTL}은 이 사실을 이용하지 않는다. 메뉴의 선마다 스크램블러 회로 전체(그림~\ref{fig:g-scr26})를
하나씩 두고 드럼 위치가 바뀔 때마다 26글자의 치환을 전부 다시 계산한다. 구조가 단순해서 맞는지 확인하기
쉽지만, 스크램블러 하나가 LUT 19,717개이므로 선이 16개인 메뉴에서는 스크램블러만으로 약 31만 5천 개다.
대상 칩(약 13만 4천 개)에 들어가지 않는다.
\item \textbf{FPGA 엔진}은 바뀌는 것만 다시 계산한다. 가운데 로터가 바뀔 때에만 안쪽 치환과 오른쪽 로터
26위치의 치환표를 여러 클럭에 걸쳐 계산해 분산 RAM에 넣어 두고, 위치마다 표에서 읽는다. 계산 회로는 한
벌만 있으면 되고 나머지는 메모리다.
\item \textbf{Rust 코어}도 같은 생각으로 짜여 있다. 가운데 로터가 바뀔 때만 표 26개를 다시 만든다
(\ref{ch:rust-bombe}장).
\end{itemize}

그림~\ref{fig:g-engine-idea}는 두 구조의 차이다. 엔진은 표를 만드는 데 클럭을 쓰지만 그것은 26위치에 한
번이다. 시간을 조금 내주고 LUT를 크게 돌려받았다. 엔진 전체는 LUT 41,040개에 들어간다. 엔진의 내부는
4부에서 다룬다.

\begin{figure}[ht]
\centering
\begin{tikzpicture}[x=1cm,y=1cm]
  \node[font=\small\bfseries] at (2.8,3.0) {기준 RTL: 위치마다 전부 다시 계산};
  \foreach \i/\yy in {0/2.1,1/1.3,2/0.5} {
    \node[blkb,minimum width=38mm,minimum height=6mm,font=\footnotesize] (s\i) at (2.2,\yy) {스크램블러 회로 (26글자분)};
    \draw[->] (s\i.east) -- ++(1.0,0);
  }
  \node at (2.2,-0.1) {$\vdots$};
  \node[note] at (2.2,-0.8) {메뉴의 선마다 한 벌 (최대 16벌)};
  \node[blkn,minimum width=10mm,minimum height=24mm] (sp) at (5.7,1.3) {전파};
  \begin{scope}[xshift=8.6cm]
  \node[font=\small\bfseries] at (3.1,3.0) {FPGA 엔진: 바뀔 때만 계산해 표에 저장};
  \node[blkb,minimum width=26mm,minimum height=10mm,font=\footnotesize] (calc) at (0.9,1.6) {치환 계산 회로\\(한 벌)};
  \node[reg,minimum width=22mm,minimum height=10mm,font=\footnotesize] (ram) at (3.9,1.6) {치환표\\(분산 RAM)};
  \node[blkn,minimum width=10mm,minimum height=24mm] (sp2) at (6.2,1.3) {전파};
  \draw[->] (calc) -- (ram); \draw[->] (ram) -- (sp2.west|-ram);
  \node[note] at (0.9,0.6) {가운데 로터가\\바뀔 때만 동작};
  \node[note] at (3.9,0.6) {위치마다 읽기};
  \end{scope}
\end{tikzpicture}
\caption{치환표를 얻는 두 가지 방법(개념도).}
\label{fig:g-engine-idea}
\end{figure}

\section{정리}

\begin{itemize}
\item FPGA의 논리는 LUT에, 기억은 플립플롭에 담긴다. 회로의 크기는 주로 LUT의 수로 잰다.
\item LUT 하나는 입력 여섯 개까지의 함수다. 출력 하나가 많은 입력에 의존할수록 LUT가 많이 든다.
\item 변수로 무엇인가를 고르는 코드는 선택 회로이고, 고를 대상의 수에 비례해 커진다.
\item 되풀이되는 계산은 표로 만들어 메모리에 두면 싸다.
\item 같은 계산을 부품을 많이 써서 빨리 할 수도, 적게 써서 천천히 할 수도 있다. 무엇이 자주 바뀌고 무엇이
드물게 바뀌는지를 보면 어디서 아낄지가 보인다.
\end{itemize}
